Our main result on Bohr sets in sumsets arising from partitions is as follows. 

%We prove \cref{th:main-kr} by adapting the proof of \cref{th:br} of Bergelson and Ruzsa. This method is quite general and so we can further attune it to obtain the following result regarding homomorphisms on compact abelian groups.

\begin{theorem} \label{th:main-partition}
Let $G$ be a compact abelian group with normalized Haar measure $\mu$ and let $\phi_1, \phi_2: G \rightarrow G$ be continuous homomorphisms satisfying
\begin{enumerate}[label=(\alph*), leftmargin=*]
\item $\phi_1, \phi_2$ are commuting, and
\item $\phi_1(G), \phi_2(G)$ have finite index in $G$.
\end{enumerate} 
Let $G = \bigcup_{i=1}^r A_i$ be a partition of $G$ into measurable sets. Then for some $1 \leq i \leq r$, 
\[
\phi_1(A_i) - \phi_1(A_i) + \phi_2(A_i)
\] 
contains a Bohr-$(k, \eta)$ set, where $k$ and $\eta$ depend only on $r$, $[G:\phi_1(G)]$ and $[G:\phi_2(G)]$.
\end{theorem}
\begin{remark}\label{remark:first_one} \
\begin{itemize}[leftmargin=*]
\item If $\mu(A) >0$, then $A + A - A$ is not guaranteed to contain a Bohr set. For a counterexample, take $G = \Z_{2N}$ for some large $N$ and $A=\{a \in \Z_{2N}: a \textup{ is odd} \}$. In particular, the analogous version of \cref{th:main-partition} for sets of positive measure fails. 

\item We do not know if the commuting condition can be removed entirely, though it can be slightly relaxed (see \cref{th:main-partition-true}). For example, commutativity is not required when $\phi_1$ or $\phi_2$ is an automorphism (see Remark \ref{rk:auto2}).
% \Hnote{Are there other interesting special cases?}
%\item \Hnote{The finite index condition should be necessary, but why????} 
%\Anote{If $\phi_i(G)$ do not have finite indices, we have a counterexample. Take $G = \T \times \Z_p$ and define $\phi_i: G \to G$ by $\phi_i(t, n) = (0, n)$. Then $\phi_1(G) + \phi_2(G) - \phi_1(G) \subseteq \{0\} \times \Z_p$ which has measure $0$ and thus cannot contain any Bohr set.}

\item The finite index condition on $\phi_1(G)$ cannot be removed, by taking for example $\phi_1 = 0$ and $\phi_2(x) = x$. On the other hand, we do not know whether the finite index condition on $\phi_2(G)$ can be removed. If we let $\phi_2 = 0$, then the situation amounts to \cref{q:kr} itself. 

%\Anote{If $\phi_i(G)$ do not have finite index, a trivial counterexample is $\phi_i(G) = 0$. In this case, $\phi_1(A_i) + \phi_2(A_i) - \phi_1(A_i) = 0$ cannot be a Bohr set. We may want a more interesting example.}
\end{itemize}
\end{remark}
